% Chapter 3 -- Data collection and field selection (draft, 17 Nov 2026)
\chapter{Data Collection and Field Selection}
\label{ch:data}
\lhead{Chapter 3. \emph{Data Collection and Field Selection}}

\noindent\emph{Status, 18 November 2026. Every data set below except the round-2 annotation and the insurance responses has been collected or built; counts are read from the repository, not from plans.}

\section{Data collected}
GRAIL uses five kinds of data: the legal text the audit is grounded in; the case records the probes are built from; authored evaluation sets; the responses of the audited models; and human labels. Table~\ref{tab:datasets} lists each with its origin, size and status.

\par\vspace{1em}
{\scriptsize
\setlength{\tabcolsep}{3pt}\setlength{\LTpre}{2.5em}\setlength{\LTpost}{1em}
\renewcommand{\arraystretch}{1.15}
\begin{longtable}{@{}>{\raggedright\arraybackslash}p{2.9cm}>{\raggedright\arraybackslash}p{3.3cm}>{\raggedright\arraybackslash}p{2.9cm}>{\raggedright\arraybackslash}p{3.1cm}>{\raggedright\arraybackslash}p{2.3cm}@{}}
\caption{Data collected for the thesis.}\label{tab:datasets}\\
\toprule
\textbf{Data} & \textbf{Origin and licence} & \textbf{Size} & \textbf{Used for} & \textbf{Status} \\
\midrule
\endfirsthead
\toprule
\textbf{Data} & \textbf{Origin and licence} & \textbf{Size} & \textbf{Used for} & \textbf{Status} \\
\midrule
\endhead
\bottomrule
\endlastfoot
EU AI Act & Official Journal L 2024/1689, English PDF from EUR-Lex \citep{euaiact2024} & 1{,}143 legal units (990 in Articles, 153 in Annexes) & Ground: retrieval index and checklist & Collected and parsed \\
Statlog German Credit & UCI repository, CC BY 4.0 \citep{statlog}; SHA-256 recorded & 1{,}000 applicants, 20 attributes, repayment label (700 repaid, 300 defaulted) & Finance stimulus (record-backed) & Collected \\
Synthetic life insurance applications & Generated from ranges declared in a stimulus pack & 12 fields per application; 14{,}825 probes & Insurance stimulus (Annex III 5(c)) & Built (18 Nov) \\
Retrieval gold sets & Authored; each gold clause id checked against the parsed Act & 10 finance + 10 insurance queries & Retrieval evaluation (E1) & Built; finance recall@3 = 1.00 with bge-base \\
Truthfulness seed bank & Authored consumer-credit questions with source hints & 20 questions & Truthfulness probes (not in the committed run) & Built \\
Finance responses & Four open models served with vLLM, temperature 0, 11 Aug--1 Sep 2026 & 61{,}259 responses, 0 errors, hash-chained log & Jury, judge, RQ1 replay & Collected \\
Judge verdicts & \texttt{microsoft/phi-4}, $k=5$, temperature 0.3 & 608 transparency items (4 models $\times$ 152) & CAL2 development & Collected \\
Human annotation, round 1 & Four raters, guideline v1.0 & 30 common items; Fleiss' $\kappa = 0.069$ & Round 1 (failed) & Collected; failed (Ch.~\ref{ch:method}) \\
Human agreement on insurance explanations & Two raters and an arbiter, guideline v2.0 & 240 items + 20 calibration & CAL1 (and CAL2) & \textbf{Pending} (runbook) \\
Insurance responses & Same four models & 14{,}825 probes per model & RQ1 confirmatory, RQ2, RQ3 (Ch.~\ref{ch:insurance}) & \textbf{Pending} (signature, amendment) \\
\end{longtable}
}

\subsection{The legal text}
The Act was downloaded as the Official Journal PDF and parsed deterministically into an Article $\to$ Paragraph $\to$ Point hierarchy: 1{,}143 units, of which 602 sit at paragraph level and 541 at point or sub-point level. Each unit keeps its citation (``Article 10(2)(f)''), heading, parent and text, so that every verdict can name the clause it rests on. The units are indexed with dense embeddings (\texttt{bge-base-en-v1.5}) and BM25. Retrieval was evaluated on ten authored queries for credit, each describing an obligation in plain words rather than quoting it; recall@3 is 1.00.

\subsection{The audited models and their responses}
Four instruction-tuned open models were audited, chosen so that one property changes at a time: Qwen2.5-7B (baseline), Llama-3.1-8B (family varies, size held), Qwen2.5-32B (size varies, family held) and Llama-3.1-70B at 4-bit (largest that fits one A6000). The list is version-controlled as provenance. Each model received the same 14{,}825 probes at temperature 0; the 7B model has additional responses from the pilot runs, which is why the total is 61{,}259 rather than $4 \times 14{,}825$. Every response is appended to a log keyed by probe hash, model and parameters, with each line chained to the previous line's hash, so a response cannot be changed or silently re-drawn.

\subsection{Human labels}
Round 1 (four raters, 30 transparency items, guideline v1.0) produced Fleiss' $\kappa = 0.069$ and is treated as a failed pilot that diagnosed the guideline. Round 2 (two raters, 300 items, guideline v1.1) is the measurement RQ1 rests on. It is not yet collected, because the guideline's two halves conflict on six of twelve items and must be reconciled and frozen first.

\section{Field selection}
Field selection happens at three levels: which clauses become requirements, which record fields the audited model sees, and which attribute is varied as the protected axis.

\subsection{Clauses}
The Act has 1{,}143 units, but only obligations whose fulfilment shows in a model's \emph{behaviour} can be probed. Five were committed and signed for credit, and the same five were derived, without change, for life insurance (Table~\ref{tab:clauses}). Obligations that concern documentation, governance or the provider's processes (for example risk-management systems or logging) are out of scope because no prompt can test them. GDPR Article 22 would be relevant to automated credit decisions but lives in a second corpus that is not loaded.

\begin{table}[H]
\centering
\vspace{1.2em}
\caption{Committed clauses and the audit dimension each becomes.}
\label{tab:clauses}
\vspace{1.2em}
\footnotesize
\begin{tabular}{@{}lll@{}}
\toprule
Clause & Obligation (short) & Dimension and probe family \\
\midrule
Art.\ 10(2)(f) & Examination for possible biases & Fairness: counterbalanced pairs \\
Art.\ 10(2)(g) & Measures to detect, prevent and mitigate biases & Fairness: counterbalanced pairs \\
Art.\ 13(1) & Transparency sufficient for deployers to interpret output & Transparency: explanation, counterfactual \\
Art.\ 15(1) & Appropriate accuracy, robustness and cybersecurity & Robustness: meaning-preserving perturbations \\
Art.\ 15(4) & Resilience to errors and inconsistencies & Robustness: meaning-preserving perturbations \\
\bottomrule
\end{tabular}
\end{table}

\subsection{Credit record fields}
Statlog has twenty attributes and a label. Eighteen attributes are rendered into each prompt, decoded from their codes into plain English (``A11'' becomes ``overdrawn''). Two are withheld and one is held out:

\begin{itemize}[noitemsep]
  \item \textbf{Personal status and sex} (attribute 9) is redacted, because it is the protected axis: the arm supplies the gender signal, and leaving the recorded value in would put two signals in one prompt.
  \item \textbf{Foreign worker} (attribute 20) is redacted for the same reason and kept as the candidate second axis for C4.
  \item \textbf{The repayment label} is never rendered. It records observed repayment, not whether an application should have been approved, so it is used as ground truth for robustness only, never as a fairness reference.
\end{itemize}

The strata (weak, marginal, strong) are terciles of a strength score built from six face-valid fields only: checking account, credit history, savings, employment duration, loan term and instalment rate. The score is not fitted to the label, since fitting it would mean building a classifier and then auditing the probes it chose, and it reads no protected field. The tercile cuts fall at scores 6 and 8 (observed range 0--15). Statlog holds no names, so a name is generated deterministically from the record identifier; it is identical in both arms of a pair and carries no information.

\subsection{Life insurance fields}
Twelve fields are sampled per application (Table~\ref{tab:lifefields}): age, sum insured, policy term, smoking status, body mass index, medical history, family history, occupation, alcohol, hazardous pursuits and a generated name. The strata are defined by age, body mass index, smoking status and medical history, with the categorical factors drawn from stratum-specific lists. Sex is never a field; it enters only through the title of the protected arm. The outcome is the quoted annual premium.

\subsection{The protected axis}
Gender is signalled by a title (Ms.\ or Mr.) on an otherwise identical application, in both sectors. A title was chosen over first names because names also signal ethnicity, age and class, so a name swap varies several attributes at once; a title varies one. It means the effect measured here is the effect of an explicit gender marker, not of a name. Every fairness pair is checked at generation to differ in the title slot and in nothing else.
